\subsection{\texorpdfstring{$p=1$}{p=1} 时的对偶}\label{sec:1.6}

对于 1-Wasserstein 度量，对偶具有十分简洁的形式. 大多数教材将这一结果视为一般代价 $c$ 的最优传输强对偶的一个特例来推导；这需要使用一些工具，将\autoref{thm:1.14}推广到二次代价 $c(x,y)=\|x-y\|^2$ 之外. 这些工具包括循环单调性和 Legendre 变换的推广概念，建议读者熟悉它们（简要介绍见\autoref{sec:1.7}）. 将这些结果用于 $p$-Wasserstein 距离，就得到以下定理；证明可参见\cite[第 3.1.1 小节]{ref281}.

\begin{theorem}\label{thm:1.20}
固定 $p\geq1$，设 $\mu,\nu\in\mathcal{P}_p(\mathbb{R}^d)$ 为两个概率测度. 则有
\begin{equation}\label{eq:1.21}
W_p^p(\mu,\nu)
=\sup_{\substack{f\in L^1(\mu),\ g\in L^1(\nu)\\
f(x)+g(y)\leq\|x-y\|^p}}
\left\{\int f\,d\mu+\int g\,d\nu\right\}.
\end{equation}
\end{theorem}

当 $p=1$ 时，可以化简这一结果，消去一个对偶函数.

\begin{theorem}\label{thm:1.21}
设 $\mu,\nu\in\mathcal{P}_1(\mathbb{R}^d)$ 为两个概率测度. 则有
\begin{equation}\label{eq:1.22}
W_1(\mu,\nu)
=\sup_{f\in\operatorname{Lip}_1}
\left\{\int f\,d\mu-\int f\,d\nu\right\},
\end{equation}
其中 $\operatorname{Lip}_1$ 是所有 1-Lipschitz 函数组成的集合.
\end{theorem}

在证明这一定理之前，先看看 Lipschitz 函数是如何出现的. 回顾\autoref{sec:1.5.2}中的半对偶，它也消去了一个对偶函数. 特别地，当 $p=1$ 时，可以将\eqref{eq:1.21}中的函数 $g$ 替换为 $c$-变换
$$
f^c(y)=\inf_{x\in\mathbb{R}^d}\{\|x-y\|-f(x)\}.
$$
有如下引理.

\begin{lemma}\label{lem:1.22}
对于 $c(x,y)=\|x-y\|$，函数 $g:\mathbb{R}^d\to\mathbb{R}$ 是某个 $c$-变换 $g=f^c$，当且仅当它是 1-Lipschitz 函数. 此外，任意 1-Lipschitz 函数都满足 $g^c=-g$.
\end{lemma}
\begin{proof}
写成
$$
g(y)=f^c(y)=\inf_{x\in\mathbb{R}^d}\{\|x-y\|-f(x)\}.
$$
对任意 $x$，由反三角不等式，函数 $y\mapsto\|x-y\|-f(x)$ 显然为 1-Lipschitz 函数. 因为 1-Lipschitz 函数的集合对取下确界封闭，所以 $g$ 也是 1-Lipschitz 函数.

为证明逆命题，设 $g$ 为 1-Lipschitz 函数，于是对任意 $x,y\in\mathbb{R}^d$，
$$
g(y)\leq\|x-y\|+g(x).
$$
对 $x$ 取下确界，得到 $g\leq(-g)^c$. 此外，
$$
(-g)^c(y)=\inf_{x\in\mathbb{R}^d}\{\|x-y\|+g(x)\}\leq g(y),
$$
其中最后一步不等式取了 $y=x$.

我们已经证明 $(-g)^c=g$，特别地，$g$ 必为一个 $c$-变换，即 $-g$ 的 $c$-变换. 引理的第二个结论由如下事实得到：若 $g$ 为 1-Lipschitz 函数，则 $-g$ 也是.
\end{proof}

现在可以证明\autoref{thm:1.21}了.

\begin{proof}[\autoref{thm:1.21}的证明]
由\autoref{sec:1.5.2}中的论证，\eqref{eq:1.21} 等于以下半对偶：
$$
\sup_{f\in L^1(\mu)}\left\{\int f\,d\mu+\int f^c\,d\nu\right\}.
$$
还可以继续优化对偶势，得到
$$
\sup_{f\in L^1(\mu)}
\left\{\int f^{cc}\,d\mu+\int(f^{cc})^c\,d\nu\right\}.
$$
由\autoref{lem:1.22}，集合
$$
\{f^{cc}:f\in L^1(\mu)\}
$$
只包含 1-Lipschitz 函数. 此外，由于 $\mu\in\mathcal{P}_1(\mathbb{R}^d)$，任意 1-Lipschitz 函数均关于 $\mu$ 可积：
$$
\begin{aligned}
\int|f|\,d\mu
&\leq\int|f(x)-f(y)|\,\mu(dx)+|f(y)|\\
&\leq\int\|x-y\|\,\mu(dx)+|f(y)|<\infty,
\end{aligned}
$$
所以 $f\in L^1(\mu)$. 因此，我们已经证明
\begin{equation}\label{eq:1.23}
\begin{aligned}
W_1(\mu,\nu)
&=\sup_{f\in\operatorname{Lip}_1}
\left\{\int f\,d\mu+\int f^c\,d\nu\right\}\\
&=\sup_{f\in\operatorname{Lip}_1}
\left\{\int f\,d\mu-\int f\,d\nu\right\},
\end{aligned}
\end{equation}
\autoref{thm:1.21}得证.
\end{proof}
